\chapter{Random matrices and the zeros of zeta}\label{ch:oc-gue}

Volume~8 has a theorem about eigenvalues: the sine kernel. Volume~14 has a theorem about zeros, on a limited range: Montgomery's pair correlation for test functions whose Fourier transform is supported in $[-1+\varepsilon,1-\varepsilon]$, under the Riemann hypothesis. They are the same formula. This chapter puts the formula on one page and says where each side stops being a theorem. Nothing proved in those two volumes is proved again. The form-factor identity, which Volume~8 cites onward to Mehta for the spacing density and does not isolate, is proved here.

\section{The form factor}

The series normalises the Fourier transform by $\hat f(\xi)=\int f(x)e^{-2\pi ix\xi}\,\dd x$. With that convention the function $\operatorname{sinc}u=\sin(\pi u)/(\pi u)$ is the transform of the indicator of $[-\tfrac12,\tfrac12]$.

\begin{theorem}[the form factor of the sine kernel]\label{thm:formfactor}
Let $K(u)=\bigl(\sin(\pi u)/(\pi u)\bigr)^2$ for $u\ne0$, and $K(0)=1$. Then $\hat K(\alpha)=(1-|\alpha|)_+$ for $|\alpha|\le1$ and $\hat K(\alpha)=0$ for $|\alpha|\ge1$. The Fourier transform of $\delta_0-K$ is $|\alpha|$ on $|\alpha|\le1$ and $1$ on $|\alpha|\ge1$.
\end{theorem}
\begin{proof}
The indicator $r=\mathbf 1_{[-1/2,1/2]}$ is even, and $\hat r=\operatorname{sinc}$. For an even function the transform is an involution, so the transform of $\operatorname{sinc}$ is $r$. The transform of a product is the convolution of the transforms, so $\hat K=r*r$. For $|\alpha|\le1$ the overlap of the two intervals has length $1-|\alpha|$, and for $|\alpha|>1$ the overlap is empty. The constant function $1$ is the transform of $\delta_0$.
\end{proof}

Volume~14 proves, under the Riemann hypothesis, that the pair correlation of the zeros has form factor $F(\alpha)=|\alpha|$ for $0\le\alpha\le1-\varepsilon$. Volume~8 proves that the sine kernel is the scaling limit of the eigenvalue process. Theorem~\ref{thm:formfactor} says that the form factor of that kernel is exactly the function Montgomery obtained on $|\alpha|\le1$, and the constant function $1$ on $|\alpha|\ge1$. The identification on $|\alpha|\le1-\varepsilon$ is a theorem, one half from each volume. The identification on $|\alpha|>1$ is a conjecture. A histogram of zeros, however far it is computed, is a computation: it does not prove the conjecture.

\begin{examplex}
Under the theorem, $\int_0^1 F(\alpha)\,\dd\alpha=\int_0^1\alpha\,\dd\alpha=\tfrac12$. Under the conjecture the integral from $0$ to $\infty$ diverges, and the meaningful comparison is the integral of $(F(\alpha)-1)$ against a test function, which on $[0,1]$ equals $\int_0^1(\alpha-1)\dd\alpha=-\tfrac12$ and on $[1,\infty)$ vanishes. The range $\alpha\le1$ is the range in which the explicit formula sees only the primes up to height about $x$: past $\alpha=1$ the prime sums are longer than the conductor, and Volume~14 stops.
\end{examplex}

\section{The story of the form factor}\label{moments:ch:zetarmt}

\begin{center}\itshape
Dyson asked what I was working on. I told him the formula, and he said: but that is the pair correlation of eigenvalues of random matrices.
\upshape(Hugh Montgomery's recollection of a conversation in 1972, in our words)
\end{center}

In 1972 Hugh Montgomery, a young number theorist, was visiting the Institute for Advanced Study in Princeton. He had
found, assuming the Riemann hypothesis, a formula for how the zeros of $\zeta$ are spaced: the density of pairs of zeros at distance $u$ (in units of
the mean spacing) seemed to be $1-\bigl(\frac{\sin\pi u}{\pi u}\bigr)^2$. At afternoon tea, the number theorist Sarvadaman Chowla insisted on
introducing him to the physicist Freeman Dyson. Montgomery hesitated to bother him; Chowla dragged him across the room. When Montgomery described his
formula, Dyson recognised it at once. It was the pair correlation of the eigenvalues of large random Hermitian matrices, which he and others had
studied as models for the energy levels of heavy nuclei.

Fifteen years later Andrew Odlyzko computed $100\,000$ zeros beyond the $10^{12}$-th, then millions near the $10^{20}$-th, and their statistics matched
the random matrix predictions to the width of the plotting pen. In 2000 Jon Keating and Nina Snaith used random matrices to guess the moments of
$|\zeta(\frac12+it)|$, a problem on which number theorists had been stuck for seventy years, and their guess reproduced the two known cases and two
deep conjectures. Nobody knows why it works. Volume~8 proves Wigner's semicircle law and the sine kernel. This chapter proves the two-by-two spacing laws, the simplest mean value theorem for $\zeta$, and explains the conjectures that connect them.

\begin{think}
Choose a large random symmetric matrix with independent standard Gaussian entries above the diagonal. Its eigenvalues are real numbers. Do you expect them
to be spread like a Gaussian? How close do you expect two neighbouring eigenvalues to come?
\end{think}

\section{The evidence}

The zeros of $\zeta$ up to height $T$ number about $\frac T{2\pi}\log\frac T{2\pi e}$ (Volume~6, the zero-counting formula), so near height $t$ their mean spacing is
$\frac{2\pi}{\log(t/2\pi)}$. To compare zeros at different heights we \emph{unfold} them: replace each ordinate $\gamma$ by
$\tilde\gamma=\frac\gamma{2\pi}\log\frac\gamma{2\pi e}+\frac78$, so that the unfolded zeros have mean spacing $1$. If the zeros were independent random points
(a Poisson process), the spacings would have density $e^{-s}$, with many near-coincidences. They do not: Figure~\ref{moments:fig:spacings} shows the unfolded
spacings of the first $5000$ zeros, which avoid $0$ and follow closely the curve $\frac{32}{\pi^2}s^2e^{-4s^2/\pi}$ that we derive in
Proposition~\ref{moments:prop:surmise} for random $2\times2$ Hermitian matrices.

\begin{center}
\refstepcounter{figure}\label{moments:fig:spacings}
\fig{moments-ch05-spacings}\\[0.5em]
{\small Figure~\thefigure. Histogram of the unfolded spacings $\tilde\gamma_{k+1}-\tilde\gamma_k$ for the first $5000$ zeros of $\zeta$, against the GUE surmise
$\frac{32}{\pi^2}s^2e^{-4s^2/\pi}$ (rust) and the Poisson density $e^{-s}$ (ochre, dashed). The zeros repel each other: small spacings are rare.}
\end{center}

The physicists' word for this is \emph{level repulsion}. For the eigenvalues of a random Hermitian matrix with Gaussian entries (the Gaussian Unitary
Ensemble, GUE), the exact spacing law of large matrices was computed by Michel Gaudin in 1961; it is very close to the $2\times2$ formula, which Eugene
Wigner had proposed in 1956 as a ``surmise''. Odlyzko's computations compare the zeros with Gaudin's exact law, and the agreement is striking. Wigner's own motivation was physics: the energy levels of a heavy nucleus are the eigenvalues of an operator far too
complicated to write down, so Wigner proposed to model it by a random matrix with the right symmetries.

\section{Level repulsion for two by two matrices}

\begin{proposition}[the Wigner surmises]\label{moments:prop:surmise}
\begin{enumerate}[label=(\alph*)]
\item Let $H=\begin{pmatrix}a&b\\b&c\end{pmatrix}$ with $a,c\sim\mathcal N(0,1)$ and $b\sim\mathcal N(0,\frac12)$ independent (the $2\times2$ GOE). The gap $D$ between its
eigenvalues, divided by its mean, has density $p_1(s)=\frac\pi2se^{-\pi s^2/4}$ on $[0,\infty[$.
\item Let $H=\begin{pmatrix}a&b\\\overline b&c\end{pmatrix}$ with $a,c\sim\mathcal N(0,1)$, and $b=u+iv$ with $u,v\sim\mathcal N(0,\frac12)$, all independent (the $2\times2$ GUE).
The normalised gap has density $p_2(s)=\frac{32}{\pi^2}s^2e^{-4s^2/\pi}$.
\end{enumerate}
\end{proposition}
\begin{proof}
The eigenvalues of $H$ are $\frac{a+c}2\pm\frac12\sqrt{(a-c)^2+4|b|^2}$, so $D=\sqrt{(a-c)^2+4|b|^2}$.

(a) $U=\frac{a-c}{\sqrt2}$ and $V=\sqrt2\,b$ are independent standard Gaussians, and $D=\sqrt2\,R$ with $R=\sqrt{U^2+V^2}$. In polar coordinates,
$\mathbb P(R\le\rho)=\frac1{2\pi}\int_0^{2\pi}\!\!\int_0^\rho e^{-r^2/2}r\dd r\dd\phi$, so $R$ has density $re^{-r^2/2}$ and mean
$\int_0^\infty r^2e^{-r^2/2}\dd r=\sqrt{\frac\pi2}$. The normalised gap $S=\frac D{\mathbb ED}=\frac R{\mathbb ER}$ has density
$\mathbb ER\cdot(\mathbb ER\,s)e^{-(\mathbb ER)^2s^2/2}=\frac\pi2se^{-\pi s^2/4}$.

(b) Now $D=\sqrt2\,R$ with $R=\sqrt{U^2+V^2+W^2}$, $U=\frac{a-c}{\sqrt2}$, $V=\sqrt2\,u$, $W=\sqrt2\,v$ independent standard Gaussians. In spherical coordinates
$R$ has density $\frac{4\pi r^2}{(2\pi)^{3/2}}e^{-r^2/2}=\sqrt{\frac2\pi}\,r^2e^{-r^2/2}$, and $\mathbb ER=\sqrt{\frac2\pi}\int_0^\infty r^3e^{-r^2/2}\dd r=2\sqrt{\frac2\pi}$. With
$\mu=\mathbb ER$, $\mu^2=\frac8\pi$, the density of $S=\frac R\mu$ is $\mu\sqrt{\frac2\pi}\mu^2s^2e^{-\mu^2s^2/2}=\frac{32}{\pi^2}s^2e^{-4s^2/\pi}$.
\end{proof}

The mechanism of repulsion is visible in the proof. A gap $D<\varepsilon$ forces two independent Gaussian quantities, or three, to be small at once,
which has probability of order $\varepsilon^2$ for real symmetric matrices and $\varepsilon^3$ for Hermitian ones. Degenerate matrices have codimension $2$
in the symmetric ones and $3$ in the Hermitian ones: this is the von Neumann--Wigner theorem of 1929, and the exponent of $s$ in the surmise, $1$ or $2$, is
that codimension minus one.

\section{Mean values of zeta}

The other side of the bridge is the size of $\zeta$ on vertical lines. To the right of the critical line its mean square is easy to guess: if
$\zeta(\sigma+it)=\sum n^{-\sigma-it}$ and the phases $n^{-it}$ behave like independent random rotations, then $|\zeta|^2$ averages to
$\sum n^{-2\sigma}=\zeta(2\sigma)$. This is a theorem for every $\sigma>\frac12$. We need two lemmas.

\begin{lemma}\label{moments:lem:zetaapprox}
For $\operatorname{Re}s=\sigma>0$, $s\ne1$, and every integer $N\ge1$,
\[
\zeta(s)=\sum_{n\le N}n^{-s}+\frac{N^{1-s}}{s-1}-s\int_N^\infty\{u\}u^{-s-1}\dd u.
\]
\end{lemma}
\begin{proof}
Let $\sigma>1$. For $M>N$, Abel summation with $\lfloor u\rfloor$ as counting function gives
$\sum_{N<n\le M}n^{-s}=M^{-s}\lfloor M\rfloor-N^{-s}N+s\int_N^M\lfloor u\rfloor u^{-s-1}\dd u$. Let $M\to\infty$ and write $\lfloor u\rfloor=u-\{u\}$:
\[
\sum_{n>N}n^{-s}=-N^{1-s}+s\int_N^\infty u^{-s}\dd u-s\int_N^\infty\{u\}u^{-s-1}\dd u=-N^{1-s}+\frac{sN^{1-s}}{s-1}-s\int_N^\infty\{u\}u^{-s-1}\dd u,
\]
and $-N^{1-s}+\frac{sN^{1-s}}{s-1}=\frac{N^{1-s}}{s-1}$. This proves the lemma for $\sigma>1$. The right side is holomorphic on $\sigma>0$, $s\ne1$, because the last
integral converges absolutely and locally uniformly there; so the identity extends by the identity theorem, $\zeta$ being continued as in Volume~5,
Chapter~10.
\end{proof}

\begin{lemma}\label{moments:lem:dirmeanvalue}
Let $\frac12<\sigma$ and $N\ge2$. Then, for every $T\ge1$,
\[
\int_T^{2T}\Bigl|\sum_{n\le N}n^{-\sigma-it}\Bigr|^2\dd t=T\sum_{n\le N}n^{-2\sigma}+O_\sigma\bigl((N^{2-2\sigma}+1)\log^2N\bigr).
\]
\end{lemma}
\begin{proof}
Expanding the square, the diagonal terms give $T\sum_nn^{-2\sigma}$, and the term of a pair $m\ne n$ is $(mn)^{-\sigma}\int_T^{2T}(n/m)^{it}\dd t$, of absolute value at
most $\frac{2(mn)^{-\sigma}}{|\log(n/m)|}$. By symmetry take $m<n$. If $m\le\frac n2$, then $\log\frac nm\ge\log2$, and these pairs contribute at most
$\frac2{\log2}\bigl(\sum_{n\le N}n^{-\sigma}\bigr)^2$. If $\frac n2<m<n$, write $m=n-r$ with $1\le r<\frac n2$: then $\log\frac nm\ge\frac{n-m}n=\frac rn$ and
$(mn)^{-\sigma}\le2^\sigma n^{-2\sigma}$, so these pairs contribute at most $2^{1+\sigma}\sum_{n\le N}n^{1-2\sigma}\sum_{r\le n}\frac1r\le2^{1+\sigma}(1+\log N)\sum_{n\le N}n^{1-2\sigma}$.
Comparing with integrals, $\sum_{n\le N}n^{-\sigma}$ and $\sum_{n\le N}n^{1-2\sigma}$ are $O_\sigma(N^{1-\sigma}+\log N)$ and $O_\sigma(N^{2-2\sigma}+\log N)$, which gives the bound.
\end{proof}

\begin{theorem}[mean square to the right of the line]\label{moments:thm:meansquare}
For every fixed $\sigma>\frac12$,
\[
\frac1T\int_T^{2T}|\zeta(\sigma+it)|^2\dd t\longrightarrow\zeta(2\sigma)\qquad(T\to\infty).
\]
\end{theorem}
\begin{proof}
Let $T\ge3$, $N=\lfloor T\rfloor$ and, for $t\in[T,2T]$ and $s=\sigma+it$, write $\zeta(s)=D(t)+R(t)+E(t)$ as in Lemma~\ref{moments:lem:zetaapprox}, with
$D(t)=\sum_{n\le N}n^{-s}$, $R(t)=\frac{N^{1-s}}{s-1}$, $E(t)=-s\int_N^\infty\{u\}u^{-s-1}\dd u$. We use the norm $\|F\|=\bigl(\int_T^{2T}|F|^2\dd t\bigr)^{1/2}$.

\emph{The main term.} By Lemma~\ref{moments:lem:dirmeanvalue}, $\|D\|^2=T\sum_{n\le N}n^{-2\sigma}+O(T^{2-2\sigma}\log^2T+\log^2T)$, and
$\sum_{n\le N}n^{-2\sigma}=\zeta(2\sigma)+O(N^{1-2\sigma})$. Since $2-2\sigma<1$ and $1-2\sigma<0$, $\|D\|^2=\zeta(2\sigma)T+o(T)$.

\emph{The term $R$.} $|R(t)|\le\frac{N^{1-\sigma}}T$, so $\|R\|^2\le\frac{N^{2-2\sigma}}T=o(T)$.

\emph{The term $E$.} With $u=e^v$, $\int_N^\infty\{u\}u^{-s-1}\dd u=\int_{\mathbb R}h(v)e^{-itv}\dd v$, where $h(v)=\{e^v\}e^{-\sigma v}$ for $v\ge\log N$ and $h(v)=0$
otherwise. The function $h$ is integrable and square integrable, so by Plancherel's theorem (Volume~8, after the change of variable $t=2\pi\xi$)
\[
\int_{\mathbb R}\Bigl|\int_{\mathbb R}h(v)e^{-itv}\dd v\Bigr|^2\dd t=2\pi\int_{\mathbb R}|h|^2\le2\pi\int_{\log N}^\infty e^{-2\sigma v}\dd v=\frac{\pi N^{-2\sigma}}\sigma.
\]
Since $|s|\le3T$ on our range, $\|E\|^2\le9T^2\frac{\pi N^{-2\sigma}}\sigma=O(T^{2-2\sigma})=o(T)$.

\emph{Conclusion.} By the triangle inequality for $\|\cdot\|$, $\bigl|\|\zeta\|-\|D\|\bigr|\le\|R\|+\|E\|=o(\sqrt T)$, and $\|D\|=\sqrt{\zeta(2\sigma)T}\,(1+o(1))$. So
$\|\zeta\|^2=\zeta(2\sigma)T(1+o(1))$.
\end{proof}

On the critical line the answer changes character. Hardy and Littlewood proved in 1918 that $\int_0^T|\zeta(\frac12+it)|^2\dd t\sim T\log T$, and Albert
Ingham in 1926 that $\int_0^T|\zeta(\frac12+it)|^4\dd t\sim\frac1{2\pi^2}T\log^4T$. No higher moment is known. Conrey and Ghosh conjectured in 1998 that
the sixth moment is $\sim42\,\frac{a_3}{9!}T\log^9T$, and Conrey and Gonek in 2001 that the eighth is $\sim24024\,\frac{a_4}{16!}T\log^{16}T$, where $a_k$ is
an explicit Euler product; the numbers $42$ and $24024$ came from long and delicate heuristics.

\begin{remark}[Keating and Snaith, 2000]\label{moments:rem:keatingsnaith}
For a random unitary matrix $U$ of size $N$, chosen with Haar measure (the circular unitary ensemble, CUE), one has the exact formula
\[
\mathbb E\,|\det(1-U)|^{2k}=\prod_{j=1}^N\frac{\Gamma(j)\Gamma(j+2k)}{\Gamma(j+k)^2}\sim\frac{G(k+1)^2}{G(2k+1)}N^{k^2},
\]
where $G$ is Barnes' function, with $G(n+1)=0!\,1!\cdots(n-1)!$. Keating and Snaith matched $N$ with $\log\frac T{2\pi}$ (both are the density of
eigenvalues, or zeros) and conjectured that
\[
\frac1T\int_0^T\Bigl|\zeta\Bigl(\frac12+it\Bigr)\Bigr|^{2k}\dd t\sim a_k\frac{G(k+1)^2}{G(2k+1)}(\log T)^{k^2}.
\]
The random matrix factor $\frac{G(k+1)^2}{G(2k+1)}$ equals $1$, $\frac1{12}$, $\frac{42}{9!}$, $\frac{24024}{16!}$ for $k=1,2,3,4$, which reproduces Hardy--Littlewood,
Ingham, Conrey--Ghosh and Conrey--Gonek; the arithmetic factor $a_k$ has to be put in by hand. The case $k=1$ of the matrix formula, $\mathbb E|\det(1-U)|^2=N+1$,
is proved in Volume~8; Montgomery's pair correlation theorem under the Riemann hypothesis is proved in Volume~14.
\end{remark}

The factor $\frac{G(k+1)^2}{G(2k+1)}$ for $k=2$ is $\frac{(0!1!)^2}{0!1!2!3!}=\frac1{12}$, and $\frac1{12}\cdot a_2$ with $a_2=\frac6{\pi^2}$ gives Ingham's
$\frac1{2\pi^2}$ (Exercise~\ref{moments:ex:barnes}).

\section{Why it is hard}

Random matrix theory predicts the local statistics of the zeros, the moments of $\zeta$, the distribution of $\log|\zeta(\frac12+it)|$, and the same
for families of $L$-functions, where the symmetry type (unitary, orthogonal, symplectic) depends on the family, as Nicholas Katz and Peter Sarnak found by
proving the analogous statements for function fields. Over $\mathbb Q$ almost none of it is proved. Montgomery's theorem determines the pair correlation only
for test functions whose Fourier transform is supported in $[-1,1]$; extending the range is equivalent to understanding correlations of primes, close to
the twin prime problem, stated in Volume~14. The moment conjectures for $k\ge3$ would need a way to control long Dirichlet polynomials on the
critical line, where the off-diagonal terms of Lemma~\ref{moments:lem:dirmeanvalue} are no longer small. And underneath, there is no matrix: the
Hilbert--P\'olya operator, if it exists, has not been found.

\section{Reading a paper: Odlyzko 1987 and Keating--Snaith 2000}

Andrew Odlyzko, \emph{On the distribution of spacings between zeros of the zeta function} (Mathematics of Computation, 1987), is a computational paper
that reads like an experimental physics report. Read the description of the computation (the zeros $10^{12}+1$ to $10^{12}+10^5$, by the
Riemann--Siegel formula), then the figures comparing the nearest-neighbour spacing and the pair correlation with the GUE predictions, and the discussion of
the small but visible deviations, which are now explained by lower-order arithmetic terms (Bogomolny and Keating). The idea worth stealing: test a
conjecture at the height where its competitors make different predictions.

Jon Keating and Nina Snaith, \emph{Random matrix theory and $\zeta(\frac12+it)$} (Communications in Mathematical Physics, 2000), computes
$\mathbb E|\det(1-U)|^{2k}$ through Selberg's integral, then deduces the moment conjecture and, as a bonus, a central limit theorem for $\log|\det(1-U)|$,
matching Selberg's theorem for $\log|\zeta|$. Read the introduction and the section with the moment formula; the idea worth stealing is to compute the
exact answer in a model where everything is computable, and to match parameters.

\section{Open doors}

\begin{itemize}
\item Compute the unfolded spacings of zeros at greater heights (Odlyzko's tables are online) and measure the deviation from the GUE law as a function of
height.
\item Compare the spacing distribution of the eigenvalues of large GOE and GUE matrices with the surmises of Proposition~\ref{moments:prop:surmise}. How large is
the difference, and where?
\item Estimate $\frac1T\int_0^T|\zeta(\frac12+it)|^6\dd t$ numerically for $T$ up to $10^4$ and compare with the Conrey--Ghosh prediction. How fast is the
convergence?
\item Do the eigenvalues of random sparse symmetric matrices, such as adjacency matrices of random regular graphs, follow the semicircle? Look up
the Kesten--McKay law.
\end{itemize}

\section{Exercises}



\begin{exercise}[$\star$]\label{moments:ex:poissonsp}
Drop $n$ independent uniform points in $[0,1]$. Show that the probability that a given point has no other point within distance $\frac sn$ to its right
tends to $e^{-s}$. This is the Poisson spacing law.
\end{exercise}

\begin{exercise}[$\star$]\label{moments:ex:surmisemean}
Check that $p_1$ and $p_2$ of Proposition~\ref{moments:prop:surmise} are probability densities with mean $1$, and that $\int_0^\varepsilon p_1\sim\frac\pi4\varepsilon^2$ and
$\int_0^\varepsilon p_2\sim\frac{32}{3\pi^2}\varepsilon^3$ as $\varepsilon\to0$.
\end{exercise}

\begin{exercise}[$\star\star$]\label{moments:ex:repulsion}
Let $(a,b,c)$ have a continuous joint density $\phi$ on $\mathbb R^3$, bounded by $M$. Show that the gap $D$ of $\begin{pmatrix}a&b\\b&c\end{pmatrix}$
satisfies $\mathbb P\bigl(D<\varepsilon,\ |a+c|\le\frac\ell2\bigr)\le\frac\pi4M\ell\varepsilon^2$ for all $\varepsilon,\ell>0$: level repulsion does not depend
on Gaussian entries.
\end{exercise}

\begin{exercise}[$\star$]\label{moments:ex:sigmagt1}
Prove Theorem~\ref{moments:thm:meansquare} for $\sigma>1$ directly from the Dirichlet series, without Lemma~\ref{moments:lem:zetaapprox}.
\end{exercise}

\begin{exercise}[$\star$]\label{moments:ex:fixedN}
Let $a_1,\dots,a_N$ be complex numbers. Show that $\frac1T\int_0^T\bigl|\sum_{n\le N}a_nn^{-it}\bigr|^2\dd t\to\sum|a_n|^2$ as $T\to\infty$.
\end{exercise}

\begin{exercise}[$\star\star$]\label{moments:ex:cue1}
For $N=1$, $U=e^{i\theta}$ with $\theta$ uniform on $[0,2\pi]$. Compute $\mathbb E|1-U|^{2k}$ directly and compare with the product formula of
Remark~\ref{moments:rem:keatingsnaith}.
\end{exercise}

\begin{exercise}[$\star$]\label{moments:ex:barnes}
Compute $\frac{G(k+1)^2}{G(2k+1)}$ for $k=1,2,3$. Given $a_2=\frac6{\pi^2}$, check that the Keating--Snaith conjecture for $k=2$ is Ingham's theorem.
\end{exercise}

\begin{exercise}[$\star\star$]\label{moments:ex:unfold}
Show that, with $\tilde\gamma=\frac\gamma{2\pi}\log\frac\gamma{2\pi e}+\frac78$, the number of unfolded zeros in $[0,X]$ is $X+O(\log X)$. What does this say about the
mean spacing of the unfolded zeros?
\end{exercise}

\section{Hints}

\begin{hint}{moments:ex:poissonsp}
The other $n-1$ points must avoid an interval of length $\frac sn$.
\end{hint}
\begin{hint}{moments:ex:surmisemean}
Use $\int_0^\infty x^{2j}e^{-x^2/2}\dd x$ and $e^{-x}=1+O(x)$.
\end{hint}
\begin{hint}{moments:ex:repulsion}
$D<\varepsilon$ means $(a-c,2b)$ lies in a disc of radius $\varepsilon$. Change variables to $(a+c,a-c,b)$.
\end{hint}
\begin{hint}{moments:ex:sigmagt1}
Expand $|\sum n^{-s}|^2$ and integrate term by term; the double series converges absolutely.
\end{hint}
\begin{hint}{moments:ex:fixedN}
$\frac1T\int_0^T(n/m)^{it}\dd t\to0$ for $m\ne n$.
\end{hint}
\begin{hint}{moments:ex:cue1}
$|1-e^{i\theta}|=2|\sin\frac\theta2|$, and $\int_0^\pi\sin^{2k}\theta\,\dd\theta$ is Wallis's integral.
\end{hint}
\begin{hint}{moments:ex:barnes}
$G(1)=G(2)=1$, $G(3)=1$, $G(4)=2$, $G(5)=12$, $G(6)=288$, $G(7)=34560$.
\end{hint}
\begin{hint}{moments:ex:unfold}
Use the formula for $N(T)$ and the monotonicity of $\gamma\mapsto\tilde\gamma$.
\end{hint}

\section{Solutions}



\begin{solution}{moments:ex:poissonsp}
\idea Independence turns ``no point in an interval'' into a power.
\writeup Condition on the given point $x$. Each of the other $n-1$ points avoids $]x,x+\frac sn]$ with probability $1-\frac sn$ (for $x\le1-\frac sn$), independently,
so the conditional probability is $(1-\frac sn)^{n-1}\to e^{-s}$. The points with $x>1-\frac sn$ have probability $\frac sn\to0$. So the probability tends to $e^{-s}$:
the normalised gap to the right neighbour is asymptotically exponential.
\end{solution}

\begin{solution}{moments:ex:surmisemean}
\idea Gaussian integrals.
\writeup With $x=\sqrt{\pi/2}\,s$,
\[
\int_0^\infty\frac\pi2se^{-\pi s^2/4}\dd s=\int_0^\infty xe^{-x^2/2}\dd x=1,
\]
and the mean is
\[
\int_0^\infty\frac\pi2s^2e^{-\pi s^2/4}\dd s=\sqrt{\frac2\pi}\int_0^\infty x^2e^{-x^2/2}\dd x=\sqrt{\frac2\pi}\sqrt{\frac\pi2}=1.
\]
For $p_2$, with $x=\sqrt{8/\pi}\,s$,
\[
\int_0^\infty p_2=\sqrt{\frac2\pi}\int_0^\infty x^2e^{-x^2/2}\dd x=1,
\]
and the mean is
\[
\sqrt{\frac\pi8}\sqrt{\frac2\pi}\int_0^\infty x^3e^{-x^2/2}\dd x=\frac12\cdot2=1.
\]
Near $0$,
$p_1(s)=\frac\pi2s(1+O(s^2))$ and $p_2(s)=\frac{32}{\pi^2}s^2(1+O(s^2))$, which integrate to $\frac\pi4\varepsilon^2$ and $\frac{32}{3\pi^2}\varepsilon^3$ at leading order.
\end{solution}

\begin{solution}{moments:ex:repulsion}
\idea The degenerate matrices form a line in $\mathbb R^3$; a tube of radius $\varepsilon$ around a segment of it has volume of order $\varepsilon^2$.
\writeup $D^2=(a-c)^2+4b^2$, so $D<\varepsilon$ means that $(a-c,2b)$ lies in the disc of radius $\varepsilon$. Change variables to $p=a+c$, $q=a-c$, $w=2b$, with Jacobian
determinant of absolute value $4$, so the density of $(p,q,w)$ is at most $\frac M4$. The event $\{D<\varepsilon,\ |p|\le\frac\ell2\}$ is the cylinder
\[
\Bigl\{|p|\le\frac\ell2,\ q^2+w^2<\varepsilon^2\Bigr\},
\]
of volume $\ell\pi\varepsilon^2$, so its probability is at most
$\frac M4\cdot\ell\pi\varepsilon^2=\frac\pi4M\ell\varepsilon^2$. For Hermitian matrices the same argument has a $3$-dimensional ball and gives
$\varepsilon^3$.
\end{solution}

\begin{solution}{moments:ex:sigmagt1}
\idea Orthogonality of $n^{-it}$ on average, with absolute convergence.
\writeup Expand as a double series:
\[
|\zeta(s)|^2=\sum_{m,n}(mn)^{-\sigma}(n/m)^{it},
\]
and $\sum_{m,n}(mn)^{-\sigma}=\zeta(\sigma)^2<\infty$, so we may integrate term by term:
\[
\frac1T\int_T^{2T}|\zeta|^2\dd t=\sum_n n^{-2\sigma}+\sum_{m\ne n}(mn)^{-\sigma}\frac1T\int_T^{2T}(n/m)^{it}\dd t.
\]
Each off-diagonal term is at most $(mn)^{-\sigma}$ and
tends to $0$, so by dominated convergence for series the second sum tends to $0$, and the limit is $\zeta(2\sigma)$.
\end{solution}

\begin{solution}{moments:ex:fixedN}
\idea A finite version of the previous exercise.
\writeup Expanding,
\[
\frac1T\int_0^T\Bigl|\sum a_nn^{-it}\Bigr|^2\dd t=\sum_n|a_n|^2+\sum_{m\ne n}a_m\overline{a_n}\frac1T\int_0^T(n/m)^{it}\dd t,
\]
and for $m\ne n$,
$\bigl|\frac1T\int_0^Te^{it\log(n/m)}\dd t\bigr|\le\frac2{T|\log(n/m)|}\to0$. There are finitely many terms.
\end{solution}

\begin{solution}{moments:ex:cue1}
\idea A Wallis integral again.
\writeup $|1-e^{i\theta}|^2=2-2\cos\theta=4\sin^2\frac\theta2$, so $\mathbb E|1-U|^{2k}=\frac1{2\pi}\int_0^{2\pi}4^k\sin^{2k}\frac\theta2\dd\theta=\frac{4^k}\pi\int_0^\pi\sin^{2k}\phi\dd\phi=\frac{4^k}\pi I_{2k}=\binom{2k}k$.
The product formula with $N=1$ gives $\frac{\Gamma(1)\Gamma(1+2k)}{\Gamma(1+k)^2}=\frac{(2k)!}{k!^2}=\binom{2k}k$. They agree.
\end{solution}

\begin{solution}{moments:ex:barnes}
\idea $G(n+1)=0!\,1!\cdots(n-1)!$.
\writeup $k=1$: $\frac{G(2)^2}{G(3)}=1$. $k=2$: $\frac{G(3)^2}{G(5)}=\frac1{12}$. $k=3$: $\frac{G(4)^2}{G(7)}=\frac4{34560}=\frac1{8640}=\frac{42}{9!}$, since $9!=362880=42\cdot8640$.
For $k=2$ the conjecture reads $\frac1T\int_0^T|\zeta|^4\sim\frac6{\pi^2}\cdot\frac1{12}\log^4T=\frac1{2\pi^2}\log^4T$, which is Ingham's theorem.
\end{solution}

\begin{solution}{moments:ex:unfold}
\idea Unfolding is the smooth part of the counting function.
\writeup By Volume~6, the zero-counting formula, $N(T)=\frac T{2\pi}\log\frac T{2\pi e}+\frac78+O(\log T)$. The map $F(\gamma)=\frac\gamma{2\pi}\log\frac\gamma{2\pi e}+\frac78$ is increasing for $\gamma\ge2\pi$,
so the number of unfolded zeros $\tilde\gamma=F(\gamma)\le X$ is $N(F^{-1}(X))$, which is $F(F^{-1}(X))+O(\log F^{-1}(X))=X+O(\log X)$. So the unfolded zeros
have density $1$: the mean spacing is $1$, up to an error that vanishes on average over long ranges.
\end{solution}

\section*{Sources and further reading}

\begin{itemize}
\item K.~Sabbagh, \emph{Dr Riemann's Zeros}, Atlantic Books, 2002, pp.~134--136, for Montgomery's account of the meeting with Dyson (the epigraph is
our paraphrase of it).
\item H.~L.~Montgomery, \emph{The pair correlation of zeros of the zeta function}, in \emph{Analytic Number Theory}, Proceedings of Symposia in Pure
Mathematics 24, AMS, 1973, 181--193.
\item A.~M.~Odlyzko, \emph{On the distribution of spacings between zeros of the zeta function}, Mathematics of Computation 48 (1987), 273--308.
\item E.~P.~Wigner, \emph{Characteristic vectors of bordered matrices with infinite dimensions}, Annals of Mathematics 62 (1955), 548--564.
M.~L.~Mehta, \emph{Random Matrices}, 3rd edition, Elsevier, 2004, for Gaudin's spacing law and the surmises.
\item G.~W.~Anderson, A.~Guionnet and O.~Zeitouni, \emph{An Introduction to Random Matrices}, Cambridge University Press, 2010, Chapter~2, for the
moment method in general.
\item G.~H.~Hardy and J.~E.~Littlewood, \emph{Contributions to the theory of the Riemann zeta-function and the theory of the distribution of primes},
Acta Mathematica 41 (1918), 119--196. A.~E.~Ingham, \emph{Mean-value theorems in the theory of the Riemann zeta-function}, Proceedings of the London
Mathematical Society (2) 27 (1926), 273--300.
\item J.~B.~Conrey and A.~Ghosh, \emph{A conjecture for the sixth power moment of the Riemann zeta-function}, International Mathematics Research
Notices 1998, 775--780. J.~B.~Conrey and S.~M.~Gonek, \emph{High moments of the Riemann zeta-function}, Duke Mathematical Journal 107 (2001), 577--604.
\item J.~P.~Keating and N.~C.~Snaith, \emph{Random matrix theory and $\zeta(1/2+it)$}, Communications in Mathematical Physics 214 (2000), 57--89.
\item N.~M.~Katz and P.~Sarnak, \emph{Zeroes of zeta functions and symmetry}, Bulletin of the AMS 36 (1999), 1--26.
\item E.~C.~Titchmarsh, \emph{The Theory of the Riemann Zeta-Function}, 2nd edition revised by D.~R.~Heath-Brown, Oxford, 1986, Chapter~7, for mean values.
\end{itemize}

